%%%%%%%%%%%%%%%%%%%%%%%%%%%%%%%%%%%%%%%%%%%%%%%%%%%%%%%%%%%%%%%%%%%%%%%%%%%%%%%
%  Laboratorio 1 - Muestreo con MCXN947
%  Estructura de carpetas esperada:
%     main.tex
%     informe-fcefyn.sty
%     img/   -> logos, capturas de Config Tools y del osciloscopio
%     src/   -> copia del codigo fuente del proyecto MCUXpresso
%
%  Todo lo que aparece entre corchetes rojos (\completar{...}) es contenido
%  que falta llenar. Para la entrega, descomentar \entregafinal: el
%  documento avisara en el log si quedo alguna marca sin resolver.
%%%%%%%%%%%%%%%%%%%%%%%%%%%%%%%%%%%%%%%%%%%%%%%%%%%%%%%%%%%%%%%%%%%%%%%%%%%%%%%
\documentclass[11pt,a4paper]{article}
\usepackage{informe-fcefyn}

% \entregafinal

% ------------------------- Datos de la caratula ------------------------------
\materia{Procesamiento Digital de Señales}
\carrera{Ingeniería en Computación}
\subtitulo{Laboratorio 1}
\tema{Muestreo con MCXN947}
\plataforma{FRDM-MCXN947 \textperiodcentered{} MCUXpresso IDE \textperiodcentered{} SDK MCUXpresso}
\autor{Costamagna, Matías}{44.512.437}{matias.costamagna@mi.unc.edu.ar}
\autor{Moreyra, Julián}{45.214.537}{julian.moreyra@mi.unc.edu.ar}
\autor{Rivera, Mariano}{}{luismarianorivera.25@mi.unc.edu.ar}
\autor{Rojas, Facundo}{}{@mi.unc.edu.ar}
% \autor{Apellido, Nombre}{Matrícula}{correo@mi.unc.edu.ar}
\docentes{Rossi, Roberto Raúl \quad\textperiodcentered\quad Parlanti, Gustavo Fernando \quad\textperiodcentered\quad Molina, Germán Alcides}
\fechaentrega{\today}
\logoizquierdo{img/logo-unc.png}
\logoderecho{img/logo-fcefyn.png}

% ---------------- Parametros del diseno (usados en las tablas) ---------------
% Frecuencia de entrada del CTIMER que dispara al ADC [Hz].
% Si se cambia el reloj en Config Tools, actualizar este valor y la tabla
% de velocidades se recalcula sola.
\frecuenciatimer{150000000}
 
\begin{document}
 
\caratula
\setcounter{tocdepth}{2}
\tableofcontents
\clearpage
 
%%%%%%%%%%%%%%%%%%%%%%%%%%%%%%%%%%%%%%%%%%%%%%%%%%%%%%%%%%%%%%%%%%%%%%%%%%%%%%%
\section{Introducción}
%%%%%%%%%%%%%%%%%%%%%%%%%%%%%%%%%%%%%%%%%%%%%%%%%%%%%%%%%%%%%%%%%%%%%%%%%%%%%%%
 
\begin{consigna}[Laboratorio 1 \textemdash{} Muestreo con MCXN947]
\textbf{Objetivo.} Manejar el módulo del conversor A/D integrado en el MCU
MCXN947.
 
\textbf{Descripción.} Realizar un programa aplicativo que sea capaz de
digitalizar una señal a través del módulo del ADC disponible en la placa
MCXN947 a distintas velocidades de muestreo; las velocidades requeridas son
\qty{8}{\kilo\sps}, \qty{16}{\kilo\sps}, \qty{22}{\kilo\sps},
\qty{44}{\kilo\sps} y \qty{48}{\kilo\sps}. Los cambios de las velocidades de
muestreo serán realizados con una de las teclas de la placa de evaluación, en
forma de un buffer circular. Cada velocidad de muestreo se indicará a través
de un color RGB del LED. Con otra tecla de la placa se habilitará la
adquisición o se parará la misma (Run/Stop). Los valores adquiridos serán
almacenados en memoria en un buffer circular de 512 muestras del tipo q15
(fraccional de 15 bits) y a su vez serán enviados a través del DAC (de
12 bits).
 
\textbf{Presentación.}
\begin{enumerate}[nosep,leftmargin=*]
  \item Mostrar el programa corriendo.
  \item Subir un archivo conteniendo lo siguiente:
  \begin{itemize}[nosep]
    \item Configuración del ambiente (pantallas con timers, ADC, etc.).
    \item Diagrama de diseño de la aplicación explicando sus partes.
    \item Código fuente de la misma con todo el proyecto.
  \end{itemize}
\end{enumerate}
\end{consigna}
 
\subsection{Requerimientos}
 
De la consigna se desprenden los siguientes requerimientos, que se usan como
guía para el diseño y como lista de verificación en la
sección~\ref{sec:verificacion}:
 
\begin{enumerate}[label=\textbf{R\arabic*},leftmargin=2.2em]
  \item \label{req:fs} Muestrear una señal analógica con el ADC a
        \qtylist{8;16;22;44;48}{\kilo\sps}.
  \item \label{req:tecla-fs} Recorrer las velocidades en forma circular con
        una tecla de la placa.
  \item \label{req:led} Indicar la velocidad activa con un color del LED RGB.
  \item \label{req:runstop} Habilitar o detener la adquisición con otra tecla
        (Run/Stop).
  \item \label{req:buffer} Guardar las muestras en un buffer circular de
        512 muestras en formato q15.
  \item \label{req:dac} Enviar cada muestra al DAC de 12 bits.
\end{enumerate}
 
\subsection{Correspondencia con la presentación}
 
\begin{table}[H]
\centering
\caption{Dónde se encuentra cada ítem pedido en la presentación.}
\label{tab:entregables}
\begin{tabular}{@{}l l@{}}
\toprule
\textbf{Ítem pedido} & \textbf{Ubicación en este informe}\\
\midrule
Configuración del ambiente & Sección~\ref{sec:ambiente}\\
Diagrama de diseño de la aplicación & Sección~\ref{sec:diseno}\\
Código fuente del proyecto & Sección~\ref{sec:implementacion} y
                             Anexo~\ref{sec:codigo}\\
Programa corriendo & Demostración en clase; evidencia en
                     sección~\ref{sec:verificacion}\\
\bottomrule
\end{tabular}
\end{table}
 
%%%%%%%%%%%%%%%%%%%%%%%%%%%%%%%%%%%%%%%%%%%%%%%%%%%%%%%%%%%%%%%%%%%%%%%%%%%%%%%
\ejercicio{Configuración del ambiente}\label{sec:ambiente}
%%%%%%%%%%%%%%%%%%%%%%%%%%%%%%%%%%%%%%%%%%%%%%%%%%%%%%%%%%%%%%%%%%%%%%%%%%%%%%%
 
\subsection{Herramientas utilizadas}
 
\begin{table}[H]
\centering
\caption{Hardware y software del entorno de trabajo.}
\label{tab:herramientas}
\begin{tabular}{@{}l l@{}}
\toprule
\textbf{Elemento} & \textbf{Versión / modelo}\\
\midrule
Placa de evaluación     & FRDM-MCXN947 (Cortex-M33, núcleo \texttt{cm33\_core0})\\
IDE                     & MCUXpresso IDE \completar{versión}\\
SDK                     & SDK\_\completar{26.6.0}\_FRDM-MCXN947\\
Herramientas de config. & MCUXpresso Config Tools (Clocks, Pins, Peripherals)\\
Proyecto base           & \completar{ejemplo del SDK usado como punto de partida}\\
Generador de señales    & \completar{marca y modelo}\\
Osciloscopio            & \completar{marca y modelo}\\
\bottomrule
\end{tabular}
